% Modernised reading — fonds Alexandre Grothendieck, shelfmark 150, pages 1-98
% (the whole folder).
%
% This is an INTERPRETATION, not a transcription. It re-reads the pages in
% today's notation and vocabulary, states the hypotheses the manuscript leaves
% implicit, and drops the critical apparatus so that the mathematics reads
% straight through. Everything the brackets and underlines carried — a doubtful
% reading, a notation he used differently, a missing passage — moves into a
% footnote. It is held to being mathematically correct as it stands; where the
% manuscript is loose or mistaken, the true statement is what appears here and
% the footnote says what the page has. In any dispute of substance the
% transcription governs: whoever cites Grothendieck must cite batch-0{1..5}.fr.tex.
%
% Scope: the folder was read whole — five batches, pages 1-98 — before any of
% this was written, and it is ONE file for the shelfmark.
%
% Notation fixed for the whole document.
%  - The page writes X^* for two different things: the open strata
%    X_i^* = X_i \ U_{j<i} X_j (pp. 2-4, 76-98), and the complement of an open
%    tube, X^* = X \ interior(T) (pp. 14-35). Here the second is written
%    X^\sharp (and X_i^\sharp = X_i ∩ X^\sharp); X_i^* keeps its first sense.
%  - Sigma^*(n) (pp. 47-63) is his raised star for spaces of strict flags; it
%    is kept, and distinct from the index * of Sigma_*, Sigma_** (pp. 40-47).
%  - « bordant » is read as « muni d'un collier » (collared), « voisinage
%    cylindrique » as « collier »; his « cône C(f) » is the mapping cylinder.
%  - \mathcal V is his round V (provoisinages), \mathbf I the interval [0,1]
%    acting multiplicatively.
%
% Substantive departures from the pages (footnoted where they occur): the
% equivalence of p. 19 needs the strict form of compatibility or the data
% X_i ∩ Y (supplied); p. 7 c) is an intersection in Sigma_2, not Sigma_12;
% p. 2 swaps id_E and id_O; p. 67 says « b fibrant » where s is meant; the
% dimension argument of p. 68 needs manifold-like Sigma_i (supplied); p. 89
% writes I' = \bar K for I' ⊂ \bar K; pp. 74-75 write X_{(d)} for X^{(d)}.
% Announced and not established: the existence of conical structures under
% non-singularity (p. 36), the independence from the filtration (p. 36), the
% formulas (72)-(73), (78) « sauf erreur » (pp. 84-86), the homotopy
% equivalences (51), (100) (« ? à vérifier », p. 96). Page 87 breaks off
% mid-sentence; page 98 ends the folder on (113).
%
% Pass: Opus 5.5 (claude-opus-5-5), 2026-10-03 — first pass, unchecked against the pages by a human.
\documentclass[11pt,a4paper]{article}
\input{../preamble/grothendieck}

\folder{150}
\pages{1}{98}
\dating{1981-1982}
\watermark{Édition de démonstration\\interprétation personnelle de l'œuvre}
\foldertitle{Espaces stratifiés et voisinages côniques (ou : déploiement des espaces stratifiés) : notes manuscrites (s.d.) — lecture modernisée du dossier entier}

\begin{document}

\section*{Résumé}

\begin{resume}
Un espace stratifié est un espace découpé en morceaux, chacun simple pris seul,
dont les complications n'apparaissent qu'aux jointures : un cône et sa pointe,
une surface et les courbes qu'elle porte, une variété algébrique et son lieu
singulier. Ces quatre-vingt-dix-huit pages, datées par lui de décembre 1981 et
janvier 1982, cherchent une manière de \emph{déplier} un tel espace : remplacer
chaque morceau par une version un peu rognée, sans plus aucune jointure, et
garder à côté la trace exacte de la façon dont les morceaux se recollaient.

L'image qui porte tout est celle d'un cône. Un voisinage d'une strate $Y$ dans
l'espace $X$ ressemble, si l'espace est raisonnable, à un cône dont la base est
le bord $\dot T$ du voisinage et dont la pointe s'étale sur $Y$ : chaque point du
bord est relié à son image dans $Y$ par un petit segment. Le dossier fait de
cette image une définition : un voisinage conique est un voisinage muni d'une
contraction $t \mapsto \lambda t$, $\lambda \in [0,1]$, qui l'écrase sur $Y$.
Ôter l'intérieur de ce voisinage, c'est découper le morceau : il reste un espace
à bord, et le bord se projette sur $Y$. Faire cela strate après strate, des plus
petites aux plus grandes, donne les \emph{strates déployées} $\Sigma_i$ — des
variétés à bord, ou à coins, quand tout est lisse — et des \emph{tubes de
raccord} $\Sigma_{ij}$, posés sur le bord de $\Sigma_j$ et fibrés au-dessus de
$\Sigma_i$. L'espace de départ se reconstruit en recollant ces pièces le long de
leurs flèches.

Le dossier avance en quatre mouvements. D'abord le vocabulaire : comment une
partition d'un espace en morceaux détermine un ordre sur les morceaux, puis les
voisinages coniques, les colliers, les compatibilités et la transversalité, et
un théorème de reconstruction par récurrence sur le nombre de strates. Ensuite
la présentation des données dépliées comme un foncteur sur l'ensemble des
drapeaux $i_0 < \cdots < i_n$ d'indices. Puis une longue digression où l'ensemble
d'indices disparaît : tout se laisse dire par un seul espace ordonné $\Sigma$,
réunion des strates déployées, dont le graphe de l'ordre est la réunion des
tubes. Enfin, en changeant de point de vue, les \emph{provoisinages} : au lieu de
choisir des voisinages, on considère le système de tous à la fois, ce qui libère
la construction des choix.

Grothendieck s'arrête au milieu de ce dernier mouvement, et il dit lui-même
où il en est. Page 55, après cinq pages de conditions : « Je m'aperçois
finalement que j'ai pas mal bafouillé pour arriver à trouver les formulations
simples — car j'ai \emph{mélangé} des propriétés purement ensemblistes …, ses
interprétations semi-simpliciales et les aspects topologiques. Je
\emph{reprends}. » La reprise tient en trois pages, et elle est nette. C'est une
leçon de méthode qui ne demande aucun bagage : séparer ce qui relève de l'ordre
seul de ce qui relève de l'espace, et les formulations se simplifient d'elles-mêmes.

Les objets que le dossier construit portent aujourd'hui plusieurs noms : voisinages en
cylindre d'application, données de contrôle de Thom et Mather, résolution ou
dépliage d'un espace stratifié en une variété à coins, catégorie des chemins de
sortie dont l'ensemble des composantes est l'ensemble ordonné des strates. Le
dossier 151, qui suit dans l'inventaire, reprend la même question par les tubes
et le recollement.

\keywords{stratified space, specialization preorder, Alexandrov topology, conical neighbourhood, mapping cylinder neighbourhood, collar, manifold with corners, Thom–Mather control data, equisingularity, unfolding of a stratified space, flag poset, internal category, ordered topological space, semi-simplicial space, pro-homotopy, Mayer–Vietoris spectral sequence}
\end{resume}

\section*{Le fil du dossier, et les conventions}

Les feuillets se suivent dans l'ordre de reliure, à deux exceptions près que
les notes signalent : les pages 29 et 30 sont interverties (la page 30 continue
la page 28, la page 31 la page 29), et les pages 92 et 93 aussi. À partir de la
page 11, une seule suite numérotée par lui, intitulée « Espaces stratifiés (étude
par voisinages coniques) », court jusqu'à la fin (ses feuillets 1 à 46) ; ses
paragraphes sont numérotés de (1) à 17.

\begin{itemize}
\item Page 1 : le plan du déploiement, avant toute définition.
\item Pages 2 à 4 : ordre de spécialisation, partitions et stratifications.
\item Pages 6, 7, 9 : le cas de trois strates, en axiomes.
\item Pages 11 à 20 : voisinages coniques, colliers, compatibilité,
équisingularité (§ 1 à 5).
\item Pages 21 à 37 : germes, transversalité, stratifications précylindriques,
et la structure conique définie par récurrence (§ 5 à 9).
\item Pages 38 à 47 : la présentation standard par un foncteur sur les drapeaux
(§ 9 et 10).
\item Pages 47 à 71 : l'espace ordonné $\Sigma$ sans ensemble d'indices, la
reprise, la récapitulation, et la catégorie topologique (§ 11 à 13).
\item Pages 72 à 75 : la construction heuristique par voisinages tubulaires (§ 14).
\item Pages 76 à 98 : provoisinages, bandes, domination (§ 15 à 17).
\end{itemize}

\emph{Conventions.} $I$ est un ensemble ordonné, fini sauf mention contraire. Un
\emph{drapeau} de $I$ est une suite strictement croissante $d = (i_0 < \cdots <
i_n)$ ; $\mathrm{Drap}(I)$ est l'ensemble des drapeaux, ordonné par inclusion.
Pour $d \subset d'$, on dit que $d$ est un \emph{segment initial} de $d'$ s'il en
est une section commençante, et qu'il est \emph{cofinal} dans $d'$ s'il a le même
plus grand élément. Les strates fermées sont les $X_i$, les strates ouvertes
$X^{*}_i = X_i \smallsetminus \bigcup_{j<i} X_j$. Le complémentaire d'un tube
ouvert est noté $X^{\sharp}$\footnote{La page écrit $X^{*}$ pour les deux objets :
les strates ouvertes (pages 2 à 4, 76 à 98) et le complémentaire $X \smallsetminus
\overset{\circ}{T}$ d'un tube (pages 14 à 35). On les distingue ici.}.

Un sous-espace fermé $Y$ de $X$ est \emph{à collier} (« partie bordante » sur la
page) s'il possède un voisinage $E$ et un homéomorphisme $E \simeq Y \times
[0,1]$ envoyant $y$ sur $(y,0)$ ; un tel $E$ est un \emph{collier} (« voisinage
cylindrique »)\footnote{Le mot « bordant » est le sien. « Collier » est le terme
reçu, au sens de Brown (1962) pour les bords de variétés ; on le prend ici pour
un fermé quelconque.}.

\pagerange{1}{1}
\section*{Le plan du déploiement (page 1)}

Le dossier s'ouvre sur ce qu'il veut obtenir\footnote{En haut à gauche de la
page 1, la date « déc. 81 -- janv. 82 », le mois de décembre lu avec doute.}. Pour
trois indices $i < j < k$, on veut des espaces $\Sigma_i, \Sigma_j, \Sigma_k$ (les
strates déployées), des tubes de raccord $\Sigma_{ij}, \Sigma_{jk}, \Sigma_{ik}$,
et une jonction $\Sigma_{ijk}$, reliés par le diagramme

\begin{tikzcd}[column sep=small, row sep=normal]
 & \Sigma_{ik} \arrow[dl] \arrow[rr, hook] & & \Sigma_{k} & \\
\Sigma_{i} & & \Sigma_{ijk} \arrow[ul, hook] \arrow[rr, hook] \arrow[dl] & & \Sigma_{jk} \arrow[ul, hook] \arrow[dl] \\
 & \Sigma_{ij} \arrow[ul] \arrow[rr, hook] & & \Sigma_{j} &
\end{tikzcd}

où les flèches $\hookrightarrow$ sont des plongements à collier (sur le bord), les
autres des fibrations propres, de fibres $F_{ij}$, $F_{jk}$, $F_{ik}$ (« fibres de
raccord »), et où les deux losanges de droite sont cartésiens. La strate
$\Sigma_i$ est « l'âme » du tube $\Sigma_{ij}$, $\Sigma_j$ en est la strate
ambiante.

Dans le cas non singulier : $\Sigma_i$ est une variété à bord dont le bord est la
réunion des $\Sigma_{ji}$, $j < i$\footnote{La page écrit $\Sigma_{ji}$ avec $j<i$,
alors que sa légende ordonne les indices croissants ; c'est bien le tube de la
strate $j$ dans la strate $i$.} ; et $\Sigma_{ik} \to \Sigma_i$ est une fibration à fibres des variétés à bord compactes, le fibré
des bords étant la réunion des $\Sigma_{ijk}$, $i<j<k$\footnote{La page écrit
« $\Sigma_{ik} \to \Sigma_k$ fibrations » ; la fibration, dans tout le reste du
dossier (pages 38, 45, 64), est la flèche vers la plus petite strate $\Sigma_i$, la
flèche vers $\Sigma_k$ étant le plongement au bord. On rétablit.}. En général, pour
un drapeau $d = (i_0 < \cdots < i_n)$ :
\begin{enumerate}
\item[a)] $\Sigma_d$ est une variété à bord, dont le bord est la réunion des
$\Sigma_{d'}$ pour $d' \supsetneq d$ et $d$ cofinal dans $d'$ ; il suffit des $d' = d
\cup \{j\}$, avec $j$ pris avant $i_0$ ou dans l'un des intervalles ouverts $]i_k,
i_{k+1}[$ ;
\item[b)] pour $\delta = (i_0, …, i_p)$ segment initial de $d$, $\Sigma_d \to
\Sigma_\delta$ est une fibration à fibres des variétés à bord compactes, le fibré
des bords étant la réunion des $\Sigma_{d'}$, $d' = d \cup \{j\}$ avec $j$ dans l'un des
intervalles $]i_k, i_{k+1}[$, $k \geq p$.
\end{enumerate}
C'est le programme du dossier ; il est repris et prouvé pas à pas aux pages 45 à
47.

\pagerange{2}{4}
\section*{L'ordre d'une partition (pages 2 à 4)}

\subsection*{Espaces et ensembles préordonnés}

Un espace topologique $X$ est préordonné par la \emph{spécialisation} : $x \leq y$
si $\overline{\{x\}} \subset \overline{\{y\}}$, c'est-à-dire si $x \in
\overline{\{y\}}$ ; c'est un ordre si et seulement si $X$ est $T_0$. Inversement un
ensemble préordonné $I$ porte la \emph{topologie d'Alexandrov} : les fermés sont
les parties stables vers le bas, les ouverts les parties stables vers le haut,
et les $I_{\geq x}$ forment une base d'ouverts ; l'adhérence de $\{x\}$ est
$I_{\leq x}$.

Une application entre ensembles préordonnés est croissante si et seulement si
elle est continue, et plus généralement une application $I \to X$ est continue
si et seulement si elle est croissante pour la spécialisation de $X$. D'où une
adjonction
\[
\mathrm{Hom}_{\mathrm{Esp}}(\mathrm{Alex}(I), X) \simeq
\mathrm{Hom}_{\mathrm{Préord}}(I, \mathrm{Spéc}(X)),
\]
avec $\mathrm{Spéc} \circ \mathrm{Alex} = \mathrm{id}_{\mathrm{Préord}}$ et une
coünité $\mathrm{Alex}(\mathrm{Spéc}\,X) \to X$ donnée par l'identité de $X$,
qui est un homéomorphisme si et seulement si toute réunion de fermés de $X$ est
fermée\footnote{La page note $\mathrm{Esp}$ le foncteur noté ici
$\mathrm{Alex}$, et $\mathrm{Ord}$ la spécialisation ; elle écrit
$\mathrm{Ord} \circ \mathrm{Esp} = \mathrm{id}_{\mathcal{E}}$ et $\mathrm{Esp}
\circ \mathrm{Ord} \to \mathrm{id}_{\mathcal{O}}$, les deux indices intervertis.
« Topologie d'Alexandrov » est le nom reçu (Alexandrov, 1937) ; la page ne le
donne pas.}. Les espaces de ce type sont ceux d'Alexandrov.

\subsection*{Pseudo-partitions et ordre récupéré}

Une application $\varphi : X \to I$ est la même chose qu'une partition de $X$ en
parties $X^{*}_i = \varphi^{-1}(i)$, éventuellement vides. Elle est continue si
et seulement si les $\bigcup_{j \geq i} X^{*}_j$ sont ouverts ; cela exige que
les $X_i = \bigcup_{j \leq i} X^{*}_j$ soient fermés, et c'est suffisant si la
famille des $X_i$ est localement finie\footnote{La page écrit, dans la formule,
la réunion des $X_j$ là où l'on attend celle des $X^{*}_j$.}.

La question posée est l'inverse : quand l'ordre de $I$ est-il \emph{déterminé} par
la partition ? Il faut que $\varphi$ soit surjective, et l'on voudrait que la
topologie de $I$ soit la topologie quotient. Pour une partition localement finie
indexée par un ensemble préordonné, avec $\varphi$ continue, on a trois
conditions :
\begin{enumerate}
\item[1°)] la topologie de $I$ est finale pour $\varphi$ : $X^{*}_J$ fermé implique
$J$ fermé ;
\item[2°)] $\varphi$ est surjective ;
\item[3°)] $\varphi$ est ouverte, ce qui revient à $\overline{X^{*}_j} =
\bigcup_{i \leq j} X^{*}_i$ pour tout $j$.
\end{enumerate}
Sous 1° et 2°, l'ordre de $I$ est l'inclusion des adhérences
$\overline{X^{*}_i}$. Une partition satisfaisant 1° à 3° équivaut à une relation
d'équivalence localement finie et ouverte, et $I$ en est l'ensemble quotient ;
$I$ est ordonné, et non seulement préordonné, si et seulement si les classes
$X^{*}_i$ sont localement fermées. On dit alors que la partition est
\emph{stratifiante}.

\subsection*{Stratifications strictes}

Du côté des strates fermées, la donnée correspondante est une famille croissante
localement finie de fermés $(X_i)_{i \in I}$, recouvrant $X$, telle que
\begin{enumerate}
\item[1°)] $i \mapsto X_i$ est croissante ;
\item[2°)] $X_i \cap X_j = \bigcup_{k \leq i,\, k \leq j} X_k$ ;
\end{enumerate}
c'est une \emph{stratification}. Elle est \emph{stricte} si de plus
\begin{enumerate}
\item[3°)] chaque $X^{*}_i = X_i \smallsetminus \bigcup_{j<i} X_j$ est non vide
(fidélité) ;
\item[4°)] chaque $X^{*}_i$ est dense dans $X_i$ (densité).
\end{enumerate}
Les partitions stratifiantes sont alors exactement les stratifications
strictes\footnote{Les conditions 3° et 4° sont des ajouts encadrés ; la seconde
est coupée par la marge et n'est lisible que par sa conclusion, qui est la
densité. Sa première formulation, « $j<i$ implique que $X_j$ rencontre $X_i$ »,
lue avec doute, n'est pas reprise.}.

Les conditions 1° et 2° passent aux images inverses, la fidélité non (un ouvert
peut éviter une strate), la densité passe aux ouverts. La fidélité au sens faible
— $X_i \subset X_j$ implique $i \leq j$ — ne suffit pas à la non-vacuité des
strates ouvertes : pour $I = \{a, b < c\}$ et $X_c = X = X_a \amalg X_b$ avec
$X_a, X_b$ non vides, l'indexation est fidèle mais $X^{*}_c$ est vide. Sous 1°,
2° et cette fidélité faible, les traces sur un ouvert forment encore une
stratification du même type, d'où un préfaisceau de stratifications et la notion
de stratification locale.

Enfin, la condition 1° se ramène, pour $I$ fini, à une condition sur un seul
indice à la fois. Si $J$ est minimal parmi les parties non fermées à $X^{*}_J$
fermé, $J$ a un plus grand élément $j$ et $J_0 = J \smallsetminus \{j\}$ est
fermé ; d'où l'équivalence de 1° avec : chaque fois que la frontière $X^{\bullet}_j
= \overline{X^{*}_j} \smallsetminus X^{*}_j$ est contenue dans $X_{J_0}$ pour un
fermé $J_0$, on a $I_{\leq j} \subset J_0 \cup \{j\}$, et la même condition vaut
dans $I_{<j}$\footnote{Le passage est encadré, en partie biffé, et le trait ne
distingue pas toujours $\bar j$ de $\bar J$ ; on a suivi la lecture que l'argument
demande, celle de la transcription.}. La page mentionne aussi la question, en
topologie modérée, d'une plus petite stratification modérée majorant une
partition donnée par des parties localement fermées modérées.

\pagerange{6}{9}
\section*{Trois strates (pages 6, 7 et 9)}

Les dessins des pages 6 et 9 montrent le cas type : un point $\Sigma_0$, une demi-droite
$\Sigma_1$ qui en part, un plan $\Sigma_2$. Les strates déployées s'obtiennent en
ôtant les tubes ouverts des strates inférieures :
\[
\Sigma_i = X_i \smallsetminus \Bigl(X_i \cap \bigcup_{k<i}
\overset{\circ}{\mathrm{Tub}}(X_k, X_i)\Bigr),
\]
et le tube de raccord $\Sigma_{ij}$ est le bord du tube de $\Sigma_i$ dans la strate
$j$ ainsi rognée, plongé dans $\Sigma_j$.

La page 7 met ce cas en axiomes. On se donne des espaces $\Sigma_0, \Sigma_1,
\Sigma_2$ et des fermés $\Sigma_{01} \subset \Sigma_0 \times \Sigma_1$, $\Sigma_{12}
\subset \Sigma_1 \times \Sigma_2$, $\Sigma_{02} \subset \Sigma_0 \times
\Sigma_2$\footnote{La page écrit « $\Sigma_{01} \subset \Sigma_{01}$ », lapsus
rétabli d'après les deux inclusions voisines.}, et l'on pose
$\Sigma_{012} = (\Sigma_{01} \times \Sigma_2) \cap (\Sigma_0 \times \Sigma_{12})$.
\begin{enumerate}
\item[a)] L'image de $\Sigma_{012}$ dans $\Sigma_0 \times \Sigma_2$ est dans
$\Sigma_{02}$ ; autrement dit, $\Sigma_0 \amalg \Sigma_1 \amalg \Sigma_2$ porte un
ordre au-dessus de $\{0 < 1 < 2\}$.
\item[b)] Les projections $\Sigma_{01} \to \Sigma_1$, $\Sigma_{02} \to \Sigma_2$,
$\Sigma_{12} \to \Sigma_2$ sont injectives, à collier ; donc $\Sigma_{012} \to
\Sigma_{12}$ aussi.
\item[c)] $\Sigma_{012} = \Sigma_{02} \cap \Sigma_{12}$, l'intersection étant prise
dans $\Sigma_2$ : le carré supérieur du diagramme est cartésien\footnote{La page
dit « dans $\Sigma_{12}$ » ; les deux termes sont des parties de $\Sigma_2$ par
b).}.
\item[d)] $\Sigma_{12} \to \Sigma_1$ est transverse au plongement à collier
$\Sigma_{01} \to \Sigma_1$ ; donc $\Sigma_{012} \to \Sigma_{12}$ est à collier.
\item[e)] $\Sigma_{02}$ et $\Sigma_{12}$ sont des parties à collier transverses de
$\Sigma_2$ ; donc $\Sigma_{012} \to \Sigma_{02}$ est à collier.
\item[f)] $\Sigma_{01} \to \Sigma_0$, $\Sigma_{02} \to \Sigma_0$, $\Sigma_{12} \to
\Sigma_1$ sont des fibrations (lisses), donc $\Sigma_{012} \to \Sigma_{02}$ et
$\Sigma_{012} \to \Sigma_{01}$ aussi.
\end{enumerate}

\pagerange{11}{20}
\section*{Voisinages coniques (pages 11 à 20)}

\subsection*{(1) La définition}

Soit $Y$ un espace. Un \emph{voisinage conique} de $Y$ est un espace $T \supset Y$,
muni d'un fermé $\dot T$ (le bord du tube) et d'une action $(\lambda, t) \mapsto
\lambda t$ du monoïde multiplicatif $\mathbf{I} = [0,1]$ telle que $\lambda y = y$
sur $Y$, $0 \cdot T \subset Y$, et que le carré

\begin{tikzcd}
\mathbf{I} \times \dot T \arrow[r] & T \\
\{0\} \times \dot T \arrow[u, hook] \arrow[r, "f"'] & Y \arrow[u, hook]
\end{tikzcd}

soit cocartésien, où $f(x) = 0 \cdot x$. Autrement dit $T$ est le \emph{cylindre de
l'application} $f : \dot T \to Y$, recollement de $\mathbf{I} \times \dot T$ et de
$Y$ le long de $\{0\} \times \dot T$, avec l'action évidente\footnote{La page dit
« le cône de l'application $f$ » et le note $\mathcal{C}(f)$. C'est le cylindre
d'application au sens de la topologie actuelle : l'extrémité $\{1\} \times \dot T$
n'est pas écrasée. Il se réduit au cône sur $\dot T$ quand $Y$ est un point, d'où le
nom. Un voisinage de cette forme est ce qu'on appelle aujourd'hui un voisinage en
cylindre d'application (« mapping cylinder neighbourhood »), notion de Siebenmann
(1972) et de Quinn ; les données de contrôle de Thom (1969) et Mather (1970) en sont la
version lisse. Le dossier ne cite aucun d'eux.}. Les $\lambda_T$ forment une
homotopie de l'identité à une rétraction de $T$ sur $Y$, fixe sur $Y$. Le bord est
déterminé par l'action : l'intérieur du tube est $\overset{\circ}{T} =
\bigcup_{\lambda < 1} \lambda T$.

Les $T_\lambda = \lambda T$ forment une famille croissante de fermés, et pour
$\lambda < 1$ l'intérieur topologique de $T_\lambda$ est $\lambda
\overset{\circ}{T}$. Pour $\varepsilon \in \,]0,1]$, l'\emph{épaississement}
$E_\varepsilon(T) = T \smallsetminus \overset{\circ}{T}_{1-\varepsilon}$ est un
voisinage de $\dot T$, homéomorphe à $[1-\varepsilon, 1] \times \dot T$.

\subsection*{(2) Voisinage conique dans un espace, colliers}

Si $Y$ est fermé dans $X$, un voisinage conique de $Y$ dans $X$ est un voisinage
$T$ muni d'une telle structure. Les $\lambda T$ en sont d'autres, et pour $\lambda
< 1$ la frontière de $\lambda T$ dans $X$ est $\lambda \dot T$ ; pour $\lambda = 1$,
ce peut être faux (prendre $X = T$).

Les colliers en sont le cas particulier où $f : \dot T \to Y$ est un
homéomorphisme. Si $T$ est un voisinage conique de $Y$, $\dot T$ est à collier
dans $T$ et dans chaque $T_\lambda$. Dans un espace compact, une réunion finie de
fermés disjoints est à collier si et seulement si chacun l'est.

\subsection*{Voisinages coniques stricts}

Soit $X^{\sharp} = X \smallsetminus \overset{\circ}{T}$, fermé de $X$.

\emph{Proposition (page 14).} Pour $\lambda \in \,]0,1[$ donné, les conditions
suivantes sont équivalentes :
\begin{enumerate}
\item[a)] $\dot T$ est à collier dans $X^{\sharp}$ ;
\item[b)] il existe un homéomorphisme de $X$ envoyant $(T, \dot T)$ sur
$(T_\lambda, \dot T_\lambda)$ ;
\item[b')] de même, avec un homéomorphisme égal à l'identité sur un $T_{\lambda'}$,
$0 < \lambda' < \lambda$ donné ;
\item[c)] il existe un voisinage conique $T'$ de $Y$ avec $T = \lambda T'$, de façon
compatible aux structures.
\end{enumerate}
Le voisinage est alors dit \emph{strict}. Sa frontière est $\dot T$, et pour tout
voisinage conique $T$ les $T_\lambda$, $\lambda < 1$, sont stricts. Une note en
marge propose de prendre les voisinages stricts comme notion de base.

\subsection*{(3) Sous-espaces coniques et compatibilité}

Une partie $T' \subset T$ est \emph{conique} si elle est stable par l'action et en
fait un voisinage conique de $Y' = T' \cap Y$\footnote{La page dit « de $Y$ » ;
c'est un voisinage conique de sa trace sur $Y$.}. Elle équivaut alors à un fermé
$\dot T' = T' \cap \dot T$ et à un fermé $Y' \supset 0 \cdot \dot T'$ de $Y$, avec
$T' = \mathbf{I} \cdot \dot T' \cup Y'$. Elle est \emph{conique stricte} si $Y' = 0
\cdot \dot T'$, c'est-à-dire si $T'$ n'a sur $Y$ que la pointe de ses propres
cônes.

Un voisinage conique $T$ de $Y$ dans $X$ est \emph{compatible} avec un fermé $X'$ si
$X' \cap T$ est conique dans $T$ ; \emph{strictement compatible} si de plus $T$ est
strict et si $\dot T$ admet dans $X^{\sharp}$ un collier $E \simeq \dot T \times
[0,1]$ dans lequel $X' \cap E$ s'écrit $X'_0 \times [0,1]$. De même, un fermé est
compatible avec un collier $E \simeq B \times [0,1]$ si sa trace y est un produit
$B' \times [0,1]$\footnote{La page remarque que la compatibilité cylindrique n'est
pas un cas particulier de la compatibilité conique, mais qu'elle lui correspond si
l'on impose la condition stricte $T' \cap Y = 0 \cdot \dot T'$, ajoutée et entourée
au bas de la page 16. En marge : « utiliser plutôt la terminologie transversale ».}.

\subsection*{(4) Équisingularité}

Un voisinage conique est \emph{équisingulier} si $f : \dot T \to Y$ est une
fibration localement triviale, et équisingulier pour une famille de parties
coniques $T_i$ si la famille des $\dot T_i \subset \dot T$ est \emph{globalement}
localement triviale au-dessus de $Y$. Si $X_i \cap Y$ est ouvert et fermé dans $Y$
connexe, on a l'alternative : ou bien $X_i \cap T = \emptyset$, ou bien $X_i
\supset Y$\footnote{L'alternative découle de la compatibilité seule : si $X_i \cap
Y$ est vide, $X_i \cap \dot T$ s'envoie dans le vide, donc est vide. La page la
présente comme conséquence de l'équisingularité, qui sert à rendre $X_i \cap Y =
0 \cdot \dot X_i$ ouvert et fermé dans le cas strict.}. Un système $(X, (X_i))$ est
équisingulier le long de $Y$ s'il existe un voisinage conique équisingulier pour
les $X_i$.

\subsection*{(5) La reconstruction}

\emph{Proposition (page 19).} Les systèmes $(X, Y, (X_i)_{i \in I}, T)$, où $T$
est un voisinage conique de $Y$ strictement compatible avec les $X_i$, équivalent
— morphismes : les isomorphismes — aux systèmes $(X^{\sharp}, \dot X, (X^{\sharp}_i),
Y, f : \dot X \to Y)$, où $\dot X$ et les $X^{\sharp}_i$ sont fermés dans
$X^{\sharp}$ ; on passe du second au premier par $X = X^{\sharp}
\amalg_{\dot X} \mathcal{C}(f)$\footnote{La page énonce l'équivalence pour $T$
seulement compatible. Elle est fausse sous cette forme : la trace $X_i \cap Y$
n'est pas déterminée par les données de droite dès qu'elle déborde l'image de
$\dot X_i$. Il faut soit la strictité au sens $X_i \cap Y = f(\dot X_i)$, que
la page impose elle-même en bas de la page 16, soit ajouter les $X_i \cap Y$ aux
données. On a retenu la première.}. Le voisinage $T$ est strict et transverse
aux $X_i$ si et seulement si $\dot X$ est à collier dans $X^{\sharp}$, compatible
avec les $X^{\sharp}_i$ ; il est équisingulier pour les $X_i$ si et seulement si
$\dot X \to Y$, avec les $\dot X_i = X_i \cap \dot X$, est une fibration
localement triviale simultanée.

La page 20, biffée, commence une définition par récurrence de la
« conification » d'une stratification finie munie d'un sous-espace à collier
$\dot X$ strictement compatible : vide si $I$ est vide, voisinage conique strict
compatible avec $\dot X$ si $I$ a un seul élément. Elle est reprise aux pages 29
à 35.

\pagerange{21}{27}
\section*{Germes et transversalité (pages 21 à 27)}

Deux voisinages coniques $T, T'$ de $Y$ dans $X$ sont \emph{équivalents} s'il existe
$\lambda, \lambda' \in \,]0,1[$ avec $\lambda T = \lambda' T'$, structures comprises
; une classe est un \emph{germe} de voisinage conique. La compatibilité avec un
fermé, la strictité, l'équisingularité et le caractère cylindrique ne dépendent
que du germe.

\emph{Équivalence de catégories (page 22).} Les systèmes $(\tilde X, \dot X,
\mathcal{E}, Y, f)$ — un espace, un fermé $\dot X$ muni d'un germe de collier
$\mathcal{E}$, un espace $Y$ et une application $f : \dot X \to Y$ — équivalent aux
systèmes $(X, Y, \mathcal{T})$ — un espace, un fermé, un germe de voisinage conique
—, par $X = \tilde X \amalg_{\dot X} Y$ : l'image du collier y devient le cylindre
de $f$. Le germe obtenu est équisingulier si et seulement si $f$ est une fibration.

\emph{Corollaire (page 23).} Les fermés $X'$ de $X$ compatibles avec le germe
conique correspondent aux couples $(\tilde X' \subset \tilde X, Y' \subset Y)$ de
fermés, $\tilde X'$ compatible avec le germe de collier, tels que $f(\dot X')
\subset Y'$ où $\dot X' = \tilde X' \cap \dot X$ ; on a $X' = \mathrm{Im}(\tilde X')
\cup Y'$, et $X'$ est conique strict si et seulement si $Y' = f(\dot X')$.

\emph{Collier compatible avec un germe conique.} Soit $X_0$ à collier dans $X$, de
germe de collier $\mathcal{E}$. On dit que $\mathcal{E}$ et le germe conique
$\mathcal{T}$ sont compatibles s'il existe des représentants $E \simeq X_0 \times
\mathbf{I}$ et $T$ tels que $T \cap E \simeq T_0 \times \mathbf{I}$, avec $\dot T
\cap E \simeq \dot T_0 \times \mathbf{I}$ et $Y \cap E \simeq Y_0 \times \mathbf{I}$
($T_0 = T \cap X_0$, etc.), la structure conique induite étant le produit par
$\mathbf{I}$ d'une structure conique de $T_0$ sur $Y_0$. Cela résume trois
conditions : les traces de $T$, $\dot T$, $Y$ sont cylindriques dans $E$ ; celles de
$E$, $\dot E$, $X_0$ sont coniques dans $T$ ; les deux familles d'homothéties
commutent.

\emph{Proposition (page 25).} Si $(X, Y, \mathcal{T})$ est décrit par $(\tilde X,
\dot X, \mathcal{E}_{\dot X}, Y, p : \dot X \to Y)$, les germes de colliers
$(X_0, \mathcal{E})$ compatibles avec $\mathcal{T}$ correspondent aux systèmes
$(\tilde X_0, \mathcal{E}_{\tilde X_0}, Y_0, \mathcal{E}_{Y_0})$ — des parties à
collier de $\tilde X$ et de $Y$ avec leurs germes — tels que : a)
$\mathcal{E}_{\tilde X_0}$ et $\mathcal{E}_{\dot X}$ sont compatibles, de sorte que
$\dot X_0 = \tilde X_0 \cap \dot X$ hérite d'un germe de collier ; b) $\dot X_0 =
p^{-1}(Y_0)$, et $p$ s'écrit $p_0 \times \mathrm{id}_{\mathbf{I}}$ sur des colliers
convenables, où $p_0 : \dot X_0 \to Y_0$\footnote{La page écrit $p_0 : X_0 \to
Y_0$, sans le point.}.

\emph{Transversalité.} Une famille finie de colliers $(B_\alpha, E_\alpha)$ est
\emph{transverse} si, pour toute partie $J$ des indices, $E_J = \bigcap_{\alpha \in J}
E_\alpha$ est isomorphe au-dessus de $B_J = \bigcap_{\alpha\in J} B_\alpha$ à $B_J
\times \mathbf{I}^J$, l'homothétie du collier $\alpha$ agissant sur la seule
coordonnée $\alpha$. C'est une structure de coin : localement, l'espace est un coin
de cube\footnote{Le nom de « variété à coins » (Cerf 1961, Douady 1961) désigne
le cas lisse de cette structure.}. La famille est transverse à un voisinage conique
$T$ de $Y$ si de plus $E_J \cap T \simeq T_J \times \mathbf{I}^J$, avec les bords et
$Y$, la structure conique étant le produit par $\mathbf{I}^J$ d'une structure conique
de $(T_J, Y_J)$.

\pagerange{27}{30}
\section*{Stratifications précylindriques (pages 27, 28 et 30)}

Soit $(X_i)_{i \in I}$ une famille croissante de fermés, $I$ fini. Elle est
\emph{précylindrique} si :
\begin{enumerate}
\item[a)] pour $j'$ prédécesseur immédiat de $j$, $X_{j'} \hookrightarrow X_j$ est
à collier ;
\item[b)] pour $j$ fixé, les colliers des prédécesseurs de $j$ sont transverses
dans $X_j$ ;
\item[c)] pour $i \leq j$, l'intervalle $[i, j]$ est un treillis booléen : si $S$
est l'ensemble des prédécesseurs immédiats de $j$ dans $[i,j]$, l'application
$k \mapsto s_k = \{ s \in S \mid s \geq k \}$ est un anti-isomorphisme de $[i,j]$ sur
$\mathfrak{P}(S)$\footnote{La lecture de b) et de la première ligne de c) est très
incertaine ; on donne la forme que la suite de la page utilise, page 28 et page
30.}.
\end{enumerate}
Ensemble, b) et c) disent que la structure locale de $X_j$ le long de $X_i$ est
cubique : un coin de cube de dimension $\mathrm{card}\, S$. Dans un treillis
booléen les bornes inférieures existent ; donc, si $j'$ est un prédécesseur de
$j$, $k \leq j$ sans être $\leq j'$, et s'il existe un minorant commun à $k$ et $j'$,
alors $k' = \inf(k, j')$ existe et est un prédécesseur de $k$. On demande alors
que le carré $(X_k, X_{k'}) \to (X_j, X_{j'})$ soit transverse : les colliers se
correspondent, $E_{k';k}$ étant l'image inverse de $E_{j';j}$\footnote{La page
écrit $E_{k;k}$ ; on attend $E_{k';k}$.}.

Un système de colliers compatible avec ces conditions fait de la stratification
une stratification \emph{cylindrique}. La page 30 ajoute enfin une condition
« oubliée », exigée désormais pour toutes les stratifications :
\[
X_i \cap X_j = \bigcup_{k \leq i,\, k \leq j} X_k ,
\]
d'où : si $X_i \cap X_j$ est non vide, $i$ et $j$ ont un minorant.

\pagerange{29}{35}
\section*{La structure conique, par récurrence (pages 29 à 35)}

Soit $(X_i)_{i \in I}$ une stratification ($X_i \cap X_j = \bigcup_{\alpha \leq i,j}
X_\alpha$, $X = \bigcup X_i$), et $(B_\alpha, \mathcal{E}_\alpha)_{\alpha \in
\Lambda}$ une famille transverse de parties à collier, transverse aux $X_i$. Soit
$I_0$ l'ensemble des éléments minimaux de $I$ ; les $X_\beta$, $\beta \in I_0$,
sont disjoints, et un germe de voisinage conique strict de $X_{I_0}$ est une famille
de germes de voisinages coniques des $X_\beta$\footnote{Les pages 29 et 30 sont
interverties dans la reliure : la page 30 continue la page 28, la page 31 la page
29.}.

Prenons $T_{I_0} = \coprod T_\beta$ strict, transverse aux colliers, et les $X_i$,
$i \in I^{*} = I \smallsetminus I_0$, strictement compatibles avec lui. Alors
$X_i \cap T_\beta = \emptyset$ si $i \not> \beta$, et $X_i \cap T_\beta =
\mathbf{I} \cdot (X_i \cap \dot T_\beta) \cup X_\beta$ si $i > \beta$\footnote{Le
premier cas demande la strictité : $0 \cdot (X_i \cap \dot T_\beta) \subset X_i
\cap X_\beta = \emptyset$ force $X_i \cap \dot T_\beta = \emptyset$.}. Le système
\[
\bigl(X, (X_i)_{i\in I}, (B_\alpha, \mathcal{E}_\alpha)_{\alpha \in \Lambda},
(T_\beta)_{\beta \in I_0}\bigr)
\]
se reconstitue donc à partir de
\[
\bigl(X^{\sharp}, (\dot X_\beta), (B^{\sharp}_\alpha, \mathcal{E}^{\sharp}_\alpha),
(X^{\sharp}_i)_{i \in I^{*}}\ ;\ (X_\beta),\ (B_{\beta\alpha},
\mathcal{E}_{\beta\alpha}),\ (p_\beta : \dot X_\beta \to X_\beta)\bigr),
\]
avec $X^{\sharp} = X \smallsetminus \overset{\circ}{T}_{I_0}$, $\dot X_\beta = \dot
T_\beta$, et les traces évidentes, par
\[
X = X^{\sharp} \amalg_{\coprod \dot X_\beta} \coprod_\beta \mathcal{C}(p_\beta),
\qquad
X_i = X^{\sharp}_i \cup \bigcup_{\beta < i} \bigl(\mathbf{I}\cdot \dot X_{i\beta}
\cup X_\beta\bigr),\quad \dot X_{i\beta} = X^{\sharp}_i \cap \dot X_\beta ,
\]
et de même pour les $B_\alpha$.

Les conditions (20) sur ces données : (a) les $X_\beta$ sont disjoints ; (b) chaque
$\dot X_\beta$ est transverse aux $(B^{\sharp}_\alpha, \mathcal{E}^{\sharp}_\alpha)$ ;
(c) les $p_\beta$ respectent les systèmes de colliers induits ; (d) les
$X^{\sharp}_i$ sont transverses à la famille formée des $\dot X_\beta$ et des
$B^{\sharp}_\alpha$. S'y ajoutent (e) $i \leq j$ implique $X^{\sharp}_i \subset
X^{\sharp}_j$, et (f) $X^{\sharp}_i \cap X^{\sharp}_j = \bigcup_{h \leq i,j}
X^{\sharp}_h$, qui disent que les $X^{\sharp}_i$ forment une stratification de
$X^{\sharp}$. La relation $X_i \cap X_j = \bigcup_{h\leq i,j} X_h$ se vérifie dans
chaque $T_\beta$ à partir de (f) sur $\dot X_\beta$, et hors des tubes par (f)
même.

\emph{Théorème (page 34).} Les systèmes du premier type, satisfaisant aux axiomes
des stratifications et aux conditions de transversalité, équivalent aux systèmes
du second, satisfaisant à (20), (e) et (f).

D'où la définition, par récurrence sur $\mathrm{card}\, I$, d'une \emph{structure
conique} sur $(X, (X_i), (B_\alpha, \mathcal{E}_\alpha))$ : vide si $I$ est vide ;
sinon, la donnée (a) d'un voisinage conique strict $T_{I_0}$ de $X_{I_0}$,
transverse aux colliers et coniquement transverse aux $X_i$, et (b) d'une structure
conique sur le système plus petit
\[
\bigl(X^{\sharp},\ (X^{\sharp}_i)_{i\in I^{*}},\ (B_\alpha, \mathcal{E}_\alpha)_{\alpha
\in \Lambda} \cup (\dot X_\beta, \mathcal{E}_\beta)_{\beta \in I_0}\bigr),
\]
où les bords des tubes rejoignent la famille des colliers. C'est, à la topologie
près, la récurrence des données de contrôle de Mather.

\pagerange{36}{37}
\subsection*{Trois remarques (pages 36 et 37)}

1) Sous des hypothèses de non-singularité, une structure conique « doit »
exister ; si la stratification est équisingulière le long de chaque strate
ouverte, il doit en exister une équisingulière, et il suffit de savoir construire
un tube équisingulier autour de $X_{I_0}$. 2) On pourrait faire la récurrence le
long de n'importe quelle filtration de $I$ par des parties dont les différences
successives sont formées d'éléments minimaux ; la notion n'en dépendrait pas
« essentiellement ». 3) La famille des $B_\alpha$ peut être remplacée par une
stratification cylindrique indexée par $\mathfrak{P}(\Lambda)^{\circ}$. Aucune
de ces trois affirmations n'est démontrée dans le dossier\footnote{Pour la
première, c'est le théorème d'existence des données de contrôle de Mather (1970)
dans le cas des stratifications de Whitney ; le dossier n'y renvoie pas.}.

Le paragraphe 9, « Description d'une stratification conique de type $I$ », ne dit
que son intention — d'abord sans stratification cylindrique supplémentaire — et la
page 37 s'arrête là.

\pagerange{38}{42}
\section*{La présentation standard (pages 38 à 42)}

Soit $\mathcal{M}$ une « catégorie topologique modèle » — les espaces topologiques,
pour fixer les idées. Une stratification équisingulière de type $I$ se présente par
un foncteur
\[
\Sigma : \mathrm{Drap}(I) \longrightarrow \mathcal{M}, \qquad (i_0, …, i_n)
\longmapsto \Sigma_{i_0 … i_n},
\]
contravariant pour l'inclusion des drapeaux, tel que :
\begin{enumerate}
\item[(1)] pour $d$ segment initial de $d'$, $\Sigma_{d'} \to \Sigma_d$ est une
fibration (il suffit de le demander pour l'oubli du dernier terme) ;
\item[(2)] l'oubli d'un terme autre que le dernier, $\Sigma_{i_0 … i_n} \to
\Sigma_{i_0 … \widehat{i_p} … i_n}$, $p \neq n$, est un plongement au bord,
à collier ;
\item[(3)] si le dernier élément $i$ de $d$ est le premier de $d'$, le carré
ci-dessous est cartésien.
\end{enumerate}

\begin{tikzcd}
\Sigma_{d \cup d'} \arrow[r, hook] \arrow[d] & \Sigma_{d'} \arrow[d] \\
\Sigma_d \arrow[r, hook] & \Sigma_i
\end{tikzcd}

Une note marginale corrige (1) : « fibrant » est trop faible pour
l'équisingularité ; il faut que, $d$ fixé, la famille des $\Sigma_{d \amalg \delta}$
au-dessus de $\Sigma_d$ soit localement triviale. Une autre ajoute la transversalité
des $\Sigma_{d'} \to \Sigma_\delta$ aux sous-espaces à collier, pour $d$ de longueur
2\footnote{La page énonce aussi un corollaire — si $d \subset d'$ ont même …,
alors $\Sigma_{d'} \to \Sigma_d$ est un plongement — dont la condition est
illisible.}.

Par (3), tout se ramène aux $\Sigma_i$, aux $\Sigma_{ij}$ ($i<j$) avec la fibration
$\Sigma_{ij} \to \Sigma_i$ et le plongement au bord $\Sigma_{ij} \to \Sigma_j$, et
aux compositions
\[
\Sigma_{ijk} = \Sigma_{ij} \times_{\Sigma_j} \Sigma_{jk} \longrightarrow \Sigma_{ik}.
\]
Les $\Sigma_i$ forment un objet $\Sigma_* = \coprod \Sigma_i$ au-dessus de $I$, les
$\Sigma_{ij}$ un objet $\Sigma_{**}$ au-dessus de $I \times I$ (avec $\Sigma_{ii} =
\Sigma_i$, et $\Sigma_{ij}$ vide si $i \not\leq j$), muni d'une source $s$ (une
fibration), d'un but $b$ (localement un plongement au bord) et d'une identité $\iota
: \Sigma_* \to \Sigma_{**}$. Avec la composition, la donnée de $\Sigma$ équivaut à
celle d'une \emph{catégorie interne} à $\mathcal{M}$, au-dessus de la catégorie
$I_{\mathcal{M}}$ définie par l'ordre de $I$. C'est même un \emph{espace ordonné} :
$\Sigma_{i_0 … i_n}$ est l'espace des chaînes $x_{i_0} \leq \cdots \leq
x_{i_n}$ de $\Sigma_*$ au-dessus de $d$, et la condition (3) se ramène au cas $n
= 2$, $\Sigma_{ijk} = \Sigma_{ik} \cap \Sigma_{jk}$ : pour $z \in \Sigma_k$,
l'ensemble $(\Sigma_*)_{\leq z}$ s'envoie injectivement dans $I_{\leq k}$, avec
l'ordre induit.

Pour reconstruire l'espace stratifié, Grothendieck considère les couples $d
\subset d'$ avec $d$ cofinal dans $d'$ ($2^n$ choix pour $d'$ de longueur $n$) ;
chacun équivaut à une chaîne $d_1 \subset \cdots \subset d_{p+1} = d'$ de segments
initiaux stricts, d'où une suite de fibrations $\Sigma_{d'} \to \Sigma_{d_p} \to
\cdots \to \Sigma_{d_1}$ et son cylindre d'application simultané $C_{\Sigma, d, d'}$.
Il reste à écrire les morphismes de transition entre ces cylindres ; la page ne
le fait pas.

\pagerange{42}{47}
\section*{Le cas des espaces topologiques (paragraphe 10, pages 42 à 47)}

Pour $\mathcal{M} = (\mathrm{Top})$, les données sont un espace $\Sigma_*$, un ordre
de graphe fermé $\Sigma_{**} \subset \Sigma_* \times \Sigma_*$, une application
croissante continue $\Sigma_* \to I$ (d'où $\Sigma_* = \coprod \Sigma_i$, $\Sigma_{**}
= \coprod_{i \leq j} \Sigma_{ij}$), avec les conditions :
\begin{enumerate}
\item[(a)] pour $i \leq j$, le but $\Sigma_{ij} \to \Sigma_j$ est un plongement
fermé ; donc, pour $z$ au-dessus de $j$ et $i \leq j$, il existe au plus un $y \leq z$
au-dessus de $i$, et $\Sigma_{ii} = \Sigma_i$ ;
\item[(b)] pour $i<j<k$, $\Sigma_{ijk} = \Sigma_{ik} \cap \Sigma_{jk}$ ;
\item[(c)] si $i, j \leq k$ n'ont pas de minorant, $\Sigma_{ik} \cap \Sigma_{jk} =
\emptyset$ ;
\item[(d)] pour $k$ fixé, les $\Sigma_{ik}$, $i<k$, sont des sous-espaces à collier
transverses de $\Sigma_k$, se coupant le long des $\Sigma_d$, $d = (i_0, …, i_n,
k)$ ;
\item[(e)] (équisingularité) les $\Sigma_{ij} \to \Sigma_i$ sont des fibrations ;
\item[(f)] (non-singularité) les $\Sigma_i$ sont des variétés à bord.
\end{enumerate}
Par (b) et (c), pour des drapeaux $d_\alpha$ de même dernier terme $k$,
$\bigcap \Sigma_{d_\alpha}$ est vide si $\bigcup d_\alpha$ n'est pas un drapeau, et
vaut $\Sigma_{\bigcup d_\alpha}$ sinon. Sous (f), les $\Sigma_{ji}$, $j<i$, forment
une stratification cylindrique du bord de $\Sigma_i$, et plus généralement les
$\Sigma_{d'}$, $d' \supsetneq d$, $d$ cofinal dans $d'$, une stratification
cylindrique du bord de $\Sigma_d$ ; l'ensemble de ces $d'$ s'identifie aux drapeaux
non vides de la somme ordinale $I_{<i_1} \amalg I_{]i_1, i_2[} \amalg \cdots
\amalg I_{]i_{n-1}, i_n[}$.

Sous (e) et (f), les fibres de $\Sigma_{ij} \to \Sigma_i$ sont des variétés à bord ;
le fibré des bords est stratifié par les $\Sigma_d$, $d$ contenant strictement
$\{i,j\}$, de premier terme $i$ et de dernier $j$, et l'ensemble de ces $d$ est
celui des drapeaux non vides de $I_{]i,j[}$. Il est vide si et seulement si $j$ est un
successeur immédiat de $i$ : c'est la condition pour que $\Sigma_{ij} \to \Sigma_i$
soit lisse, c'est-à-dire à fibres sans bord. Plus généralement, pour $d =
(i_1, …, i_p)$ segment initial strict de $d' = (i_1, …, i_n)$, le fibré des
bords de $\Sigma_{d'} \to \Sigma_d$ est stratifié par les drapeaux non vides de
$I_{]i_p, i_{p+1}[} \amalg \cdots \amalg I_{]i_{n-1}, i_n[}$, et la fibration est
lisse si et seulement si chacun de $i_p, i_{p+1}, …, i_n$ est successeur
immédiat du précédent — « à partir de $i_p$, pas de $i_1$ ! ». C'est l'énoncé b) de
la page 1, désormais démontré.

\pagerange{47}{54}
\section*{Digression : l'espace ordonné sans ensemble d'indices (paragraphe 11, pages 47 à 54)}

L'ensemble $I$ disparaît. On se donne un espace ordonné $\Sigma$ : un espace
$\Sigma(0)$ et un fermé $\Sigma(1) \subset \Sigma(0) \times \Sigma(0)$, graphe d'un
ordre. Soit $\Sigma^{*}(1)$ le graphe de l'ordre strict et $\Sigma^{*}(n)$ l'espace
des chaînes strictes $x_0 < \cdots < x_n$ ; ils forment un espace semi-simplicial
(sans dégénérescences)\footnote{La page dit « pré-simplicial », et note ces
espaces d'une petite étoile, gardée ici.}. Les conditions :
\begin{enumerate}
\item[(a)] le but $\Sigma(1) \to \Sigma(0)$ est une immersion propre localement à
collier ; la diagonale est alors ouverte et fermée dans $\Sigma(1)$ ;
\item[(b)] pour tout carré cocartésien d'ensembles totalement ordonnés $\delta = d
\cap d'$, $d'' = d \cup d'$, avec $\delta$ segment initial de $d$ contenant le
dernier élément de $d$, le carré des $\Sigma^{*}$ correspondant est cartésien ; et
pour $d = (i_0, …, i_n)$, au voisinage de chaque point, $\Sigma^{*}(d)$ est
l'intersection transverse des sous-espaces à collier $\Sigma^{*}(i_\alpha, i_n) \to
\Sigma^{*}(i_n)$ ; de même pour toute inclusion cofinale ;
\item[(c)] (équisingularité) la source $\Sigma^{*}(1) \to \Sigma^{*}(0)$ est propre et
fibrante ; donc $\Sigma^{*}(d') \to \Sigma^{*}(d)$ l'est pour $d$ segment initial de
$d'$ ;
\item[(d)] (non-singularité) $\Sigma^{*}(0)$ est une variété à bord dont le bord est
l'image du but ;
\item[(e)] (mixte) les fibres de la source sont des variétés à bord, le fibré des
bords étant l'image de $\Sigma^{*}(2) \to \Sigma^{*}(1)$, $(x<y<z) \mapsto (x<z)$.
\end{enumerate}
Une marge doute : le bord de $\Sigma^{*}(d)$ est-il bien la réunion des
$\Sigma^{*}(d')$, $d \subset d'$ cofinal ? « Si c'est vrai, ça risque d'être
délicat ! »

La condition (b) est ensuite précisée en trois points. Pour $x \in \Sigma$ :
b$_1$) les $z_\alpha \in \Sigma^{*}(1)$ de but $x$ correspondent bijectivement aux
germes en $x$ de sous-espaces à collier $B_\alpha$, images des germes de
$\Sigma^{*}(1)$ en les $z_\alpha$ ; b$_2$) si les $z_\alpha$ sont au nombre de $n$, il
existe un unique $\tilde x \in \Sigma^{*}(n)$ dont ils sont les images par les $n$
faces ; b$_3$) pour $1 \leq d \leq n$, les points de $\Sigma^{*}(d)$ au-dessus de $x$
correspondent aux parties à $d$ éléments de $\{z_\alpha\}$, et le germe de
$\Sigma^{*}(d)$ en un tel point a pour image l'intersection des $B_\alpha$
correspondants.

D'où une reformulation. Les inclusions cofinales d'ensembles totalement ordonnés
non vides forment une catégorie équivalente à celle des ensembles totalement
ordonnés $d$, éventuellement vides, par $d \mapsto d \amalg \{\omega\}$. Le foncteur
$\Gamma(d) = \Sigma^{*}(d \amalg \{\omega\})$ est donc un espace semi-simplicial
au-dessus de $\Gamma(\emptyset) = \Sigma$, et b$_1$ à b$_3$ disent qu'il provient
d'une \emph{stratification cylindrique locale ordonnée} de $\Sigma$ : les strates de
codimension 1 passant par $x$ sont totalement ordonnées, continûment en $x$, et
$\Gamma(d)$ est l'espace des germes de strates de codimension $d+1$.

Autre description : $\Gamma'(d)$, l'espace des suites de $d$ germes distincts de
strates de codimension 1 en un point, est la puissance relative $d$-ième de $\Gamma_0
= \Sigma^{*}(1)$ sur $\Sigma$ privée des diagonales, et $\Gamma'(d) \simeq
\mathfrak{S}_d \times \Gamma(d)$. Un ordre total relatif est la donnée d'une partie
ouverte et fermée $\Gamma_1$ de $\Gamma_0 \times_\Sigma \Gamma_0$ privé de la
diagonale, telle que $\Gamma_1$ et son symétrique la partagent, et transitive :
$(\Gamma_1 \times_\Sigma \Gamma_0) \cap (\Gamma_0 \times_\Sigma \Gamma_1)$ s'envoie
par $\mathrm{pr}_{13}$ dans $\Gamma_1$.

\pagerange{55}{63}
\section*{La reprise (pages 55 à 63)}

« Je m'aperçois finalement que j'ai pas mal bafouillé … Je
\emph{reprends}. » Soit $\Sigma$ un ensemble ordonné muni d'une topologie séparée,
$\Sigma(1)$ le graphe de son ordre.
\begin{enumerate}
\item[(a)] Le but $b = \mathrm{pr}_2 : \Sigma(1) \to \Sigma$ est propre et
localement une immersion à collier. Alors la diagonale est ouverte et fermée
dans $\Sigma(1)$, son complémentaire est $\Sigma^{*}(1)$, et $\Sigma(1)$ est fermé
dans $\Sigma \times \Sigma$.
\item[(b)] Pour tout $x$, l'ensemble fini $\Sigma_{\leq x}$ est totalement ordonné.
Les chaînes de longueur $d$ de but $x$ correspondent donc aux parties à $d+1$
éléments de $\Sigma_{\leq x}$, et la fibre de $\Sigma^{*}(d) \to \Sigma$ en $x$ à
$\mathfrak{P}_d(\{1, …, N(x)\})$, $N(x) = \mathrm{card}\, \Sigma_{<x}$\footnote{La
page écrit $N(x) = \mathrm{card}(\Sigma_{\leq x})$ ; pour que les chaînes strictes
de longueur $d$ finissant en $x$ correspondent aux parties à $d$ éléments, il faut
compter $\Sigma_{<x}$.}. Tout cela est purement ensembliste.
\item[(c)] Pour tout $x$, les germes $B_\alpha$ de sous-espaces à collier
correspondant aux $z_\alpha$ de but $x$ sont deux à deux distincts et mutuellement
transverses.
\end{enumerate}
Les germes d'intersection des $B_\alpha$ correspondent alors aux points des
$\Sigma^{*}(d)$ au-dessus de $x$ : $\Sigma$ est muni d'une stratification
cylindrique locale ordonnée, $\Sigma^{*}(1)$ étant l'espace des strates de
codimension 1.

\subsection*{Retrouver l'ordre}

Inversement, partons de la stratification et de l'ordre relatif. Pour retrouver
l'ordre de $\Sigma$, on se donne une application source $s : \Sigma^{*}(1) \to
\Sigma$, injective sur chaque fibre du but et d'image ne rencontrant pas la
diagonale — c'est-à-dire $(s, b) : \Sigma^{*}(1) \to \Sigma \times \Sigma$ injective,
d'image disjointe de la diagonale —, et l'on demande qu'il existe $p_2 :
\Sigma^{*}(2) \to \Sigma^{*}(1)$ avec $b \circ p_2 = s \circ b'$ et $s \circ p_2 = s
\circ s'$ ; $p_2$ est alors unique, et l'on exige que le carré

\begin{tikzcd}
\Sigma^{*}(2) \arrow[r, hook, "b'"] \arrow[d, "p_2"'] & \Sigma^{*}(1) \arrow[d, "s"] \\
\Sigma^{*}(1) \arrow[r, hook, "b"'] & \Sigma
\end{tikzcd}

soit cartésien, et que $s$ soit transverse à l'immersion à collier $b$.

Posons $x \prec y$ s'il existe $u \in \Sigma^{*}(1)$ avec $s(u) = x$, $b(u) = y$.
C'est transitif : pour $u : x \prec y$ et $v : y \prec z$, on a $b(u) = s(v)$, donc
par le carré cartésien un unique $\delta \in \Sigma^{*}(2)$ avec $b'(\delta) = v$,
$p_2(\delta) = u$, et $w = s'(\delta)$ vérifie $s(w) = s(p_2 \delta) = x$, $b(w) =
b(b'\delta) = z$\footnote{La page écrit $b(s'(d)) = b(b'(d))$ : c'est l'identité
simpliciale $b \circ s' = b \circ b'$ des deux faces qui conservent le dernier
sommet.}. C'est irréflexif, puisque l'image de $(s, b)$ évite la diagonale. Donc
$\prec$ est l'ordre strict d'un ordre sur $\Sigma$, dont $\Sigma^{*}(1)$ est le graphe
strict et $\Sigma^{*}(2)$ l'espace des chaînes $x \prec y \prec z$, les faces étant
$(x,y,z) \mapsto (x,y), (x,z), (y,z)$.

\emph{Scholie (pages 60 à 62).} La situation est décrite par : 1°) un espace
$\Sigma$ muni d'une stratification précylindrique, dont l'espace des strates de
codimension 1 s'immerge par $b$, immersion propre à collier ; 2°) un ordre total
relatif sur ces strates, donné par une partie fermée de $\Sigma^{*}(1) \times_\Sigma
\Sigma^{*}(1)$ ; 3°) une application source $s$ telle que a) $(s, b)$ est injective
et d'image disjointe de la diagonale, b) $s \circ s'$ se factorise par $(s, b)$ en un
$p_2$ rendant le carré ci-dessus cartésien, avec $s$ transverse à $b$.

Ce qui se redit, en termes d'espace ordonné : a) chaque $\Sigma_{<x}$ est totalement
ordonné ; b) le but $b$ est une immersion propre, donc les $\Sigma_{<x}$ sont finis ;
c) les germes en $x$ images des $(\Sigma^{*}(1), z_\alpha)$ forment une famille
transverse de sous-espaces à collier ; d) $s$ est transverse à $b$. Puis trois cas :
équisingulier ($s$ fibrante propre), non singulier ($\Sigma$ variété à bord de bord
$b(\Sigma^{*}(1))$), et les deux à la fois ($s$ fibration propre à fibres des variétés
à bord, le fibré des bords étant l'image de $s' : \Sigma^{*}(2) \to \Sigma^{*}(1)$).
Une marge suggère que le troisième est conséquence des deux premiers.

Les flèches principales forment l'hexagone

\begin{tikzcd}[column sep=small]
 & \Sigma^{*}(1) \arrow[r, hook, "b"] \arrow[dl, "s"'] & \Sigma & \\
\Sigma & \Sigma^{*}(2) \arrow[u, hook, "s'"] \arrow[rr, hook, "b'"] \arrow[d, "p_2"'] & & \Sigma^{*}(1) \arrow[ul, hook, "b"'] \arrow[dl, "s"] \\
 & \Sigma^{*}(1) \arrow[r, hook, "b"] \arrow[ul, "s"] & \Sigma &
\end{tikzcd}

aux deux carrés cartésiens, où $b, b', s'$ sont des immersions à collier propres et
$s, p_2$ des fibrations propres. Ces fibrations ne sont pas surjectives : l'image de
$s$ est l'ensemble des points non maximaux, celle de $p_2$ celui des $x<y$ avec $y$
non maximal.

\pagerange{64}{66}
\section*{Récapitulation (paragraphe 12, pages 64 à 66)}

Les données de déploiement sur un ensemble ordonné fini $I$ sont des espaces
$\Sigma_i$, des espaces $\Sigma_{ij}$ ($i<j$), des applications $s_{ij} : \Sigma_{ij}
\to \Sigma_i$ et $b_{ij} : \Sigma_{ij} \to \Sigma_j$, avec :
\begin{enumerate}
\item[a)] $b_{ij}$ est un plongement fermé, par lequel on identifie $\Sigma_{ij}$ à
une partie de $\Sigma_j$ ;
\item[b)] $\Sigma_{ik} \cap \Sigma_{jk} = \emptyset$ si $i, j<k$ sont incomparables ;
\item[c)] pour $i<j<k$, on pose $\Sigma_{ijk} = \Sigma_{ik} \cap \Sigma_{jk} \subset
\Sigma_k$, et l'on a $\Sigma_{ijk} = s_{jk}^{-1}(\Sigma_{ij})$ dans
$\Sigma_{jk}$\footnote{La page écrit $b_{ij}^{-1}(\Sigma_{ij})$ ; dans $\Sigma_{jk}$,
c'est l'image inverse par la fibration $s_{jk}$ vers $\Sigma_j$ que l'hexagone
demande.} ;
\item[d)] les $\Sigma_{ji} \hookrightarrow \Sigma_i$ sont à collier, et en $x \in
\Sigma_i$ ceux qui contiennent $x$ — totalement ordonnés par b) — sont mutuellement
transverses\footnote{La page écrit « les $\Sigma_j \hookrightarrow \Sigma_{ji}$ ».} ;
\item[e)] $s_{ij}$ est transverse au système des $b_{ki}$.
\end{enumerate}
D'où l'hexagone de la page 64, semblable au précédent, pour $i<j<k$. Avec $\Sigma =
\coprod \Sigma_i$ et $\Sigma^{*}(1) = \coprod_{i<j} \Sigma_{ij}$, les conditions a) à c)
disent que $(s, b)$ est injective, d'image le graphe d'un ordre strict, que $b$ est
une immersion propre, que $\varphi : \Sigma \to I$ est strictement croissante et
injective sur chaque $\Sigma_{<x}$, lequel est totalement ordonné ; d) est la
condition c) de la scholie. Cas équisingulier : les $s_{ij}$ fibrations propres ; non
singulier : les $\Sigma_i$ variétés à bord, de bord la réunion des $\Sigma_{ji}$ ;
les deux : $s_{ij}$ fibration propre à fibres des variétés à bord, le fibré des bords
étant la réunion des $\Sigma_{ijk}$.

\pagerange{66}{71}
\section*{Une digression : catégories topologiques (paragraphe 13, pages 66 à 71)}

Soit $\underline{\Sigma}$ une catégorie topologique — objets $\Sigma(0)$, flèches
$\Sigma(1)$, source et but $s, b$, composition — dont la source est une fibration.
Alors $\pi_0(\Sigma(0))$ est un ensemble préordonné, le graphe étant l'image de
$\pi_0(\Sigma(1))$ dans $\pi_0(\Sigma(0)) \times \pi_0(\Sigma(0))$. La relation est
réflexive par les identités, transitive parce que $\pi_0(\Sigma(2)) \to
\pi_0(\Sigma(1)) \times_{\pi_0 \Sigma(0)} \pi_0(\Sigma(1))$ est surjective, par le
lemme suivant appliqué à la source\footnote{La page écrit « grâce à l'hypothèse
que $b$ est fibrant » ; l'hypothèse posée à la page 66, et celle dont le lemme a
besoin pour $\Sigma(2) = \Sigma(1) \times_{b, \Sigma(0), s} \Sigma(1)$, porte sur
$s$.}.

\emph{Lemme (page 67).} Dans un carré cartésien d'espaces $X' = X \times_Y Y'$, avec
$f : X \to Y$ une fibration, $\pi_0(X') \to \pi_0(X) \times_{\pi_0(Y)} \pi_0(Y')$ est
surjective\footnote{Ici $\pi_0$ désigne les composantes connexes par arcs, et
« fibration » une fibration de Hurewicz : on relève un chemin qui joint $f(x)$ à
l'image de $y'$. Pour les composantes connexes, il faudrait supposer les espaces
localement connexes par arcs.}.

Le foncteur canonique $\underline{\Sigma} \to I = \pi_0(\Sigma(0))$, vers la
catégorie discrète de l'ensemble préordonné, découpe $\Sigma(0) = \coprod \Sigma_i$ et
$\Sigma(1) = \coprod_{i \leq j} \Sigma_{ij}$ ; il est universel parmi les foncteurs
vers un ensemble préordonné. Supposons de plus :
\begin{enumerate}
\item[a)] le but $\Sigma(1) \to \Sigma(0)$ est une immersion propre, de sorte que
l'identité $\delta : \Sigma(0) \to \Sigma(1)$ est un plongement ouvert et fermé et
$\Sigma(1) \simeq \delta(\Sigma(0)) \amalg \Sigma^{*}(1)$ ;
\item[b)] la restriction du but à $\Sigma^{*}(1)$ est une immersion d'images
localement rares ;
\item[c)] les composantes de $\Sigma(0)$ sont de dimension finie.
\end{enumerate}
Alors $\Sigma^{*}_{ii} = \Sigma_{ii} \cap \Sigma^{*}(1)$ est vide : fibré sur
$\Sigma_i$, il serait de dimension au moins $\dim \Sigma_i$, et, d'image rare par une
immersion, de dimension moindre. Donc $\Sigma_i \simeq \Sigma_{ii}$ : toute flèche
entre deux objets d'une même composante est une identité\footnote{L'argument de
dimension demande plus que la page ne dit : que la fibration soit localement
triviale, pour que la dimension totale majore celle de la base, et que les $\Sigma_i$
soient de dimension pure et à la manière des variétés, pour qu'une partie fermée
rare soit de dimension strictement moindre. C'est le cas dans le cadre non singulier
qui est celui du dossier ; on suppose ces hypothèses.}. Si de plus d) toute flèche
inversible est une identité — par exemple si $\underline{\Sigma}$ vient d'un espace
ordonné —, $I$ est \emph{ordonné} et $\Sigma \to I$ est strictement croissante.

\emph{Conséquence (page 70).} Les déploiements équisinguliers de stratifications
indexés par $I$, au sens de la récapitulation, équivalent aux données : (a) d'un
espace ordonné $\Sigma$ satisfaisant aux conditions a), b), c) et au cas
équisingulier de la scholie, dont les composantes connexes sont de dimension finie,
tel qu'en posant $\tilde I = \pi_0(\Sigma)$, chaque $b_{ij} : \Sigma_{ij} \to
\Sigma_j$, $i<j$ dans $\tilde I$, soit un plongement fermé ; (b) d'une application
strictement croissante, non nécessairement surjective, $\tilde I \to I$. La condition
b) de la récapitulation se démontre : un point de $\Sigma_{ik} \cap \Sigma_{jk}$
donnerait deux éléments $< x$ au-dessus de $i$ et de $j$, comparables puisque
$\Sigma_{<x}$ est totalement ordonné, donc $i$ et $j$ comparables.

\emph{Remarque (page 71).} Deux contextes de « déploiement » ont produit de tels
espaces ordonnés — $s$ fibration propre, $b$ immersion propre, deux points distincts
d'une même composante incomparables. Pour une stratification par les $X_i$, $\Sigma =
\coprod X_i$ et $\Sigma(1) = \coprod_{X_i \subset X_j} X_i$ : $s$ est étale, rien de
spécial sur $b$. Pour la théorie présente, $b$ est à collier et rien de spécial sur
$s$. Dans les deux cas, passer à $\tilde I = \pi_0(\Sigma)$ revient à remplacer les
strates par leurs composantes connexes\footnote{Le premier contexte est, en termes
actuels, la catégorie des chemins de sortie d'une stratification, dont l'ensemble
des composantes est l'ensemble ordonné des strates (MacPherson ; Treumann 2009 ;
Lurie, \emph{Higher Algebra}, annexe A). Le dossier précède ces travaux.}.

\pagerange{72}{75}
\section*{Construction heuristique des $\Sigma_i$, $\Sigma_{ij}$ (paragraphe 14, pages 72 à 75)}

À préciser « considérablement » en topologie modérée. Pour $\Sigma_i$, on ôte de
$X_i$ l'intérieur d'un voisinage tubulaire standard $T_{<i}$ de $\bigcup_{k<i} X_k$ :
$\Sigma_i = X_i \smallsetminus \overset{\circ}{T}_{<i}$. Pour $\Sigma_{ij}$, on prend un
voisinage tubulaire $T'_{ij}$ de $\Sigma_i$ dans $X_j$, petit devant $T_{<i}$, son bord
$\dot T'_{ij}$, la trace $S \dot T'_{ij}$ des $X_k$, $i<k<j$ — vide si $j$ est
successeur immédiat de $i$ —, et l'on ôte un voisinage tubulaire de cette partie
singulière : $\Sigma_{ij} = \dot T'_{ij} \smallsetminus \overset{\circ}{T}(S \dot
T'_{ij}) \subset \Sigma_j$.

Pour que les $\Sigma_{ik}$ et $\Sigma_{jk}$ soient en position correcte dans
$\Sigma_k$, on fixe une fonction dimension strictement croissante $\delta : I \to
\mathbb{N}$, on pose $X_{(d)} = \bigcup_{\delta(i) \leq d} X_i$, et l'on construit une
filtration décroissante par des fermés $X = X^{(0)} \supset X^{(1)} \supset \cdots
\supset X^{(N+1)} = \emptyset$ : $T_d$ voisinage tubulaire de $X_{(d)}$, petit devant
les précédents, et $X^{(d)} = X \smallsetminus \bigcup_{\alpha < d}
\overset{\circ}{T}_\alpha$. On pose alors, pour $\delta(i) = d$,
\[
\Sigma_i = X_i \cap X^{(d)} = X_i \smallsetminus \bigcup_{j<i} X_i \cap
\overset{\circ}{T}_j, \qquad \Sigma_{ij} = \dot T_i \cap \Sigma_j ,
\]
\[
\Sigma_{i_0 … i_n} = \dot T_{i_0} \cap \cdots \cap \dot T_{i_{n-1}} \cap
\Sigma_{i_n} = \Sigma_{i_0 i_n} \cap \cdots \cap \Sigma_{i_{n-1} i_n}.
\]
La page écrit ici $X_i \cap X_{(d)}$\footnote{Pour $X_i \cap X^{(d)}$ ; elle numérote
aussi (49) deux formules différentes.}.
Les morphismes de transition pour les inclusions cofinales de drapeaux sont
évidents ; les autres dépendent d'un choix de rétractions des $T_i^{(d)}$ sur
$X_i^{(d)} = \Sigma_i$, et pour que les $\Sigma_{ij} \to \Sigma_i$ soient des
fibrations il faut des hypothèses d'équisingularité et la locale trivialité de ces
rétractions.

\pagerange{76}{83}
\section*{Provoisinages (paragraphe 15, pages 76 à 83)}

Le point de vue change : au lieu de choisir un voisinage, on prend le système de
tous. Soit $\mathcal{V}_i$ le \emph{provoisinage} tubulaire de la strate ouverte
$X^{*}_i$ dans $X$ — le système projectif de ses voisinages, vu comme pro-espace —,
et $\mathcal{V}_{ij} = \mathcal{V}_i \cap \mathcal{V}_j$\footnote{« Provoisinage »
est lu de l'interligne « pro » des pages 76 et 79. La notion de pro-espace est
celle de la théorie de la forme et de l'homotopie propre (Artin--Mazur 1969 ;
Edwards--Hastings 1976) ; l'« entrelacs homotopique » de Quinn (1988) en est une
version.}. L'inclusion $X^{*}_i \to \mathcal{V}_i$ est supposée une équivalence
d'homotopie de pro-espaces\footnote{Ce n'est pas vrai pour une stratification
quelconque ; c'est vrai si les strates ouvertes sont des rétractes par déformation
de voisinages, par exemple sous l'hypothèse conique du début du dossier. La page
96 met elle-même un « ? à vérifier ! » en face d'un énoncé de ce type.}, et
$\mathcal{V}_{ij}$ est vide si $i, j$ sont incomparables dès que $X$ est
héréditairement normal, par exemple métrisable : $X^{*}_i$ et $X^{*}_j$ sont alors
séparés, aucune n'adhérant à l'autre, donc à voisinages disjoints\footnote{La page
dit « si les ouverts de $X$ sont normaux ». La justification : $X^{*}_i \cap
\overline{X^{*}_j} \subset X^{*}_i \cap X_j$, vide si $i \not\leq j$.}. L'espace
s'obtient par recollement :
\[
X = \varinjlim \bigl(\mathcal{V}_i, \mathcal{V}_{ij} \rightrightarrows\bigr) =
\varinjlim_{d \in \mathrm{Drap}(I)} \mathcal{V}_d, \qquad \mathcal{V}_{i_0 … i_n}
= \mathcal{V}_{i_0} \cap \cdots \cap \mathcal{V}_{i_n},
\]
limite prise niveau par niveau sur des voisinages ouverts recouvrant $X$.

Relativement à une partie localement fermée $Z$ : $\mathcal{V}^Z_d = \mathcal{V}_d
\cap Z$, et $\mathcal{V}^Z_{i_n} \to \mathcal{V}_{i_n}$ est une équivalence d'homotopie
si $Z$ contient $X^{*}_{i_n}$. Pour une partie $J$ de $I$, $X^{*}_J = \bigcup_{j \in J}
X^{*}_j$ ; c'est un fermé si $J$ est un \emph{idéal} (stable vers le bas), et une
partie localement fermée si $J$ est une \emph{bande} (convexe : $i<k<j$, $i,j \in J$
impliquent $k \in J$), car $X^{*}_J = X_{\tilde J} \smallsetminus X_{J'}$ où $\tilde J$
est l'idéal engendré et $J' = \tilde J \smallsetminus J$\footnote{Ici $\tilde J$ est
l'idéal engendré ; à partir de la page 84, la page note $\bar J$ l'idéal engendré
et $\tilde J$ le filtre engendré. On suit chaque fois sa notation, en le
rappelant.}. Les $X_i \cap X^{*}_J$ stratifient $X^{*}_J$, et
\[
X^{*}_J \simeq \varinjlim_{d \in \mathrm{Drap}(J)} \mathcal{V}^{J}_d .
\]

Pour deux bandes $J, L$, on s'intéresse à $\mathcal{V}^L_J = \mathcal{V}_J \cap X^{*}_L$,
« provoisinage de second ordre ». Il ne change pas si l'on remplace $J$ par ses
éléments qui minorent un élément de $L$, et $L$ par ses éléments qui majorent un
élément de $J$ ; on se ramène ainsi au cas $J \prec L$ : tout élément de $J$ est
majoré par un élément de $L$, et tout élément de $L$ majore un élément de $J$.

\emph{Lemme (page 81).} Pour deux parties $J, L$ de $I$, sont équivalents : a) $J$ et
$L$ sont des bandes et $J \prec L$ ; b) il existe une bande $I'$ dont $J$ est un idéal
et $L$ un filtre, chaque élément de $I'$ majorant un élément de $J$ et minorant un
élément de $L$. La bande $I'$ est alors unique : $I' = \{ i \mid \exists j \in J, l \in
L,\ j \leq i \leq l \}$, la plus petite bande contenant $J \cup L$\footnote{La fin de
la démonstration dit « un élément de $I$ » pour « de $J$ ».}.

Alors $\mathcal{V}^L_J$ se calcule dans $X' = X^{*}_{I'}$, où $X'_J$ est fermé et
$X'_L$ ouvert :
\[
\mathcal{V}^L_J = \bigcup_{\substack{j \in J,\ l \in L\\ j \leq l}}
\mathcal{V}^{X'}_{j,l} = \varinjlim_{d \in \mathrm{Drap}^L_J(I)} \mathcal{V}^{X'}_d ,
\]
où $\mathrm{Drap}^L_J(I)$ est l'ensemble des drapeaux de $J \cup L$ commençant dans $J$
et finissant dans $L$ ; les intersections deux à deux sont vides sauf si la réunion
des indices est un drapeau, d'au plus quatre termes. Pour un faisceau $F$ sur
$\mathcal{V}^L_J$, d'où une suite spectrale
\[
E_2^{pq} = H^p\bigl(\mathrm{Drap}^L_J(I),\ d \mapsto H^q(\mathcal{V}^{X'}_d, F)\bigr)
\Longrightarrow H^{p+q}(\mathcal{V}^L_J, F).
\]
Les $H^p$ sont ceux de l'ensemble ordonné à coefficients dans le
foncteur\footnote{C'est-à-dire les foncteurs dérivés de la limite. C'est la suite
spectrale d'un recouvrement ouvert indexé par un ensemble ordonné, stable par
intersection (l'intersection de deux membres est le membre de la réunion des
drapeaux, ou vide), à la manière de Mayer--Vietoris.}.

\pagerange{84}{87}
\section*{$\Sigma_J$ et $\Sigma_{J,L}$ (paragraphe 16, pages 84 à 87)}

On revient aux strates déployées, désormais indexées par des bandes. Notons $\bar
J$ l'idéal engendré par $J$ et $\tilde J$ le filtre engendré. Pour une bande $J$,
\[
\Sigma_J = X^{*}_J \smallsetminus X^{*}_J \cap \bigcup_{i \in \bar J \smallsetminus J}
\overset{\circ}{T}_i ,
\]
de sorte que $\Sigma_\emptyset = \emptyset$, $\Sigma_{\{j\}} = \Sigma_j$, $\Sigma_I =
X$, et $\Sigma_J = X_J$ si $J$ est un idéal. « Sauf erreur »,
\[
\Sigma_J \simeq \varinjlim_{d \in \mathrm{Drap}(J)} \Sigma_d, \qquad \Sigma_J
\hookrightarrow X^{*}_J \ \text{équivalence d'homotopie},
\]
et le bord $\dot \Sigma_J$ est la réunion des images des $\Sigma_{ij} \to \Sigma_j$,
$i<j$ dans $J$, limite des $\Sigma_d$ pour $d$ d'au moins deux termes.

Pour deux bandes disjointes $J \prec L$, avec $I' = \tilde J \cap \bar L$ et $J' = I'
\smallsetminus L$, on a des idéaux $J \subset J' \subset I'$ et des inclusions $\Sigma_J
\subset \Sigma_{J'} \subset \Sigma_{I'}$. Le tube de $\Sigma_J$ dans $\Sigma_{I'}$
est $\Sigma_{I'} \cap \bigcup_{j \in J} T_j$, et l'on pose
\[
\Sigma_{J,L} = \dot T_{\Sigma_J, \Sigma_{I'}} \cap \Sigma_L \ \simeq\
\varinjlim \Sigma_d ,
\]
la limite étant prise sur les drapeaux de $J \cup L$ de premier terme dans $J$ et de
dernier dans $L$ (« sauf erreur ») ; chacun s'écrit uniquement $(j_1 … j_p\, l_1
… l_q)$, $p, q \geq 1$, d'où l'inclusion $\Sigma_{J,L} \hookrightarrow \Sigma_L$
et la projection $\Sigma_{J,L} \to \Sigma_J$, induites par $\Sigma_{j_1 … j_p l_1
… l_q} \hookrightarrow \Sigma_{l_1 … l_q}$ et $\Sigma_{j_1 … j_p l_1
… l_q} \to \Sigma_{j_1 … j_p}$.

\emph{Remarque (page 87).} Les $\Sigma_{J,L}$ et les inclusions $\Sigma_{J,L} \to
\Sigma_L$ ont un sens, dans une catégorie isotopique, indépendamment de tout choix
de rétractions coniques locales autour des strates, du moins en topologie modérée ;
la projection $\Sigma_{J,L} \to \Sigma_J$, elle, utilise ces rétractions. La phrase
suivante s'arrête sur « il est plus commode, » ; la page 88 ouvre un nouveau
paragraphe.

\pagerange{88}{98}
\section*{Récapitulation sur les provoisinages (paragraphe 17, pages 88 à 98)}

\subsection*{Topologie inférieure, bandes, domination}

Les parties fermées de $I$ — stables vers le bas — sont les fermés de la topologie
\emph{inférieure}, celle d'Alexandrov des pages 2 à 4 : l'adhérence de $A$ est $\bar
A = \{ i \mid \exists j \in A,\ i \leq j \}$, et le plus petit ouvert contenant $B$ est
$\tilde B = \{ i \mid \exists j \in B,\ i \geq j \}$. Une partie est localement fermée
si et seulement si $J = \bar J \cap \tilde J$, si et seulement si c'est une bande.
Toute intersection de bandes est une bande, et la bande engendrée par $J$ est $B(J) =
\bar J \cap \tilde J$.

La \emph{domination} $J \prec K$, définie par $J \subset \bar K$ et $K \subset \tilde J$,
est un préordre sur les parties de $I$ ; $J \sim K$ si et seulement si $B(J) = B(K)$, et
l'ensemble ordonné quotient est celui des bandes, ordonné par domination. Pour deux
bandes $J \prec K$, la bande engendrée par $J \cup K$ est $I' = \tilde J \cap \bar
K$ ; $J$ y est fermé et codense ($I' \subset \tilde J$), $K$ ouvert et dense ($I'
\subset \bar K$)\footnote{La page écrit $I' = \bar K$ ; c'est une inclusion, la
densité de $K$ dans $I'$.}, et réciproquement.

\subsection*{Stratifications ordinaires}

Une \emph{stratification ordinaire} de type $I$ est une application $i \mapsto X_i$
satisfaisant (Strat 1) les $X_i$ sont fermés et forment un recouvrement localement
fini ; (Strat 2) $i \mapsto X_i$ est croissante ; (Strat 3) $X_i \cap X_j =
\bigcup_{k \leq i,\, k \leq j} X_k$\footnote{La page annonce « les deux conditions » et
en énumère trois.}. Pour un drapeau $d$, posons $X_d = X_{i_0}$ : on obtient un
espace semi-simplicial au-dessus de $X$, résolution de $X$, d'où pour un faisceau
abélien $F$
\[
E_2^{pq} = H^p\bigl(I,\ i \mapsto H^q(X_i, F)\bigr) = H^p_{ss}\bigl(\mathrm{Drap}(I),\
d \mapsto H^q(X_d, F)\bigr) \Longrightarrow H^{p+q}(X, F),
\]
suite spectrale d'un recouvrement fermé localement fini\footnote{Au sens de
Godement, \emph{Topologie algébrique et théorie des faisceaux}, II.5.2 ; la page ne
le cite pas.}, « souvent assez peu pratique » parce que les $X_i$ sont compliqués.

Les strates ouvertes $X^{*}_i = X_i \smallsetminus \bigcup_{j<i} X_j$ non vides forment
une partition de $X$. Axiomes supplémentaires, à n'invoquer que par mention
expresse : (Strat 4) fidélité, $X^{*}_i \neq \emptyset$ — non stable par restriction à
un ouvert ; (Strat 4 bis) les $X^{*}_i$ sont connexes non vides ; (Strat 5) densité,
$X_i = \overline{X^{*}_i}$, c'est-à-dire $X_j$ rare dans $X_i$ pour $j<i$ — stable
par restriction à un ouvert\footnote{La ligne de (Strat 4 bis) finit sur « i.e. » et
la suivante porte (Strat 5) ; on lit deux axiomes distincts.}. Moyennant la densité,
la fidélité équivaut à $X_i \neq \emptyset$.

On a $X_i = \bigcup_{j \leq i} X^{*}_j$. Si la stratification est fidèle, $X_i \subset
X_j$, $X^{*}_i \subset X_j$ et $i \leq j$ sont équivalents : $I$ s'identifie à une
famille de parties de $X$, et une stratification fidèle équivaut à une famille
localement finie de fermés non vides, recouvrant $X$, satisfaisant (Strat 3), dont
aucun membre n'est réunion des membres strictement plus petits\footnote{Les pages
92 et 93 sont interverties dans la reliure : la page 91 se poursuit par la page 93,
puis par la page 92.}.

Pour $J \subset I$, $X^{*}_J = \bigcup_{j \in J} X^{*}_j$ et $X_J = \bigcup_{j \in J}
X_j = X^{*}_{\bar J}$. L'application $J \mapsto X^{*}_J$ commute aux réunions et aux
intersections quelconques. On a $X^{*}_J = X^{*}_{\bar J} \smallsetminus X^{*}_{\bar J
\smallsetminus J}$ ; si $J$ est une bande, $\bar J \smallsetminus J$ est fermé et $X^{*}_J$
localement fermé ; sous la densité, $\overline{X^{*}_J} = X_J = X^{*}_J \amalg
X^{*}_{\bar J \smallsetminus J}$. Si $L$ est ouvert, $X^{*}_L$ est le complémentaire de
$X^{*}_{I \smallsetminus L}$, donc ouvert. Sous densité et fidélité, les réciproques
valent : $X^{*}_J$ est fermé si et seulement si $J$ est fermé, $X^{*}_L$ ouvert si et
seulement si $L$ l'est, $X^{*}_J$ localement fermé si et seulement si $J$ est une bande.

\subsection*{Les provoisinages, définitivement}

Soit $\mathcal{V}_i = \mathcal{V}(X^{*}_i, X)$ le pro-espace des voisinages de
$X^{*}_i$. Alors
\[
\mathcal{V}_i \cap X^{*}_j \neq \emptyset \iff X^{*}_i \cap \overline{X^{*}_j} \neq
\emptyset \Longrightarrow X^{*}_i \cap X_j \neq \emptyset \Longrightarrow i \leq j ,
\]
la première implication étant une équivalence sous la densité, la seconde sous la
fidélité ; pour une stratification dense et fidèle, $\mathcal{V}_i \cap X^{*}_j \neq
\emptyset$ si et seulement si $i \leq j$. La relation n'est pas symétrique ; en
revanche, si $X$ est héréditairement normal, $\mathcal{V}_i \cap \mathcal{V}_j \neq
\emptyset$ implique que $i$ et $j$ sont comparables, avec réciproque sous la
fidélité\footnote{« Non vide » s'entend du pro-espace : tout voisinage rencontre.
$\mathcal{V}_i \cap X^{*}_j$ est non vide en ce sens si et seulement si $X^{*}_i$
rencontre l'adhérence de $X^{*}_j$, sinon le complémentaire de cette adhérence est
un voisinage qui l'évite.}.

On pose $\mathcal{V}_{i_0 … i_n} = \mathcal{V}_{i_0} \cap \cdots \cap
\mathcal{V}_{i_n}$ pour les drapeaux ; ce sont les « pièces de construction »
élémentaires des types d'homotopie, par limites inductives finies si $I$ est fini.
Pour $Z$ localement fermé contenant $X^{*}_{i_n}$, $\mathcal{V}^Z_{i_0 … i_n} =
\mathcal{V}_{i_0 … i_n} \cap Z \to \mathcal{V}_{i_0 … i_n}$ « est » une
équivalence d'homotopie — la marge porte « ? à vérifier ! ». Pour une bande $J$, et
pour deux bandes $J, L$ :
\[
\mathcal{V}_J = \mathcal{V}(X^{*}_J, X) = \bigcup_{j \in J} \mathcal{V}_j =
\varinjlim_{d \in \mathrm{Drap}(J)} \mathcal{V}_d , \qquad X^{*}_J \hookrightarrow
\mathcal{V}^Z_J \hookrightarrow \mathcal{V}_J \ \ (Z \supset X_J),
\]
\[
\mathcal{V}_{J,L} = \mathcal{V}_J \cap \mathcal{V}_L = \varinjlim_{\substack{d \in
\mathrm{Drap}(J \cup L)\\ d \cap J \neq \emptyset,\ d \cap L \neq \emptyset}}
\mathcal{V}_d \quad (J \prec L), \qquad \mathcal{V}^L_J = \mathcal{V}_J \cap X^{*}_L
\hookrightarrow \mathcal{V}_{J,L} ,
\]
cette dernière inclusion étant une équivalence d'homotopie si $J \prec L$ sont des
bandes. Dans le cas $J \prec L$, les drapeaux en jeu commencent dans $J$, finissent
dans $L$, et leurs termes dans $J$ forment un segment initial, ceux dans $L$ un
segment final. Une marge propose la notation $\mathcal{V}^{\flat}_{i_0 … i_n} =
\mathcal{V}_{i_0} \cap \cdots \cap \mathcal{V}_{i_{n-1}} \cap X^{*}_{i_n}$.

Enfin, pour une chaîne de bandes $J_0 \prec J_1 \prec \cdots \prec J_n$,
\[
\mathcal{V}_{J_0, …, J_n} = \mathcal{V}_{J_0} \cap \cdots \cap \mathcal{V}_{J_n} =
\bigcup_{j_\alpha \in J_\alpha} \mathcal{V}_{j_0 … j_n} = \varinjlim_{\substack{d
\in \mathrm{Drap}(J_0 \cup \cdots \cup J_n)\\ d \cap J_\alpha \neq \emptyset\ \forall
\alpha}} \mathcal{V}_d .
\]
Le dossier s'arrête sur cette formule ; un bloc de calculs raturés au bas de la
dernière page n'a pas pu être lu.

\end{document}
